\documentclass[12pt]{interim_report}
\usepackage{amsmath,amsthm,graphicx}

\title{Modelling terahertz pulse generation due to ultrafast photocurrents in semiconductors }
\author{Yaoqi Dong, Yuhao Chang, Shenglong Li}
\date{12 December 2025}
\supervisor{Professor Oleg Mitrofanov }
\aiuse{{\bf If you have not used generative AI make sure that {\tt\bf\textbackslash aiuse} is empty. Only edit and include text inside {\tt\bf\textbackslash aiuse} if you have used generative AI.}

I acknowledge the use of name and version of generative AI system  Chat-GPT-  (Publisher OpenAI, URL of the AI system  \\https://chat.openai.com) to brief description of context in which AI tool was used to summarise my initial notes, help with LaTeX format and proofread my final draft.}
\wordcount{4401}

\begin{document}

\coverpage
\declaration
\maketitle

\tableofcontents

%% ============================================================
%% EXECUTIVE SUMMARY
%% ============================================================
\section{Executive Summary}
% TODO: Write a concise executive summary (no more than one page).
% - Briefly what is this project about?
% - Discuss the main points -- one paragraph each
% - Include the main results and findings
% - Spark the reader's interest!

%% ============================================================
%% INTRODUCTION AND LITERATURE REVIEW
%% ============================================================
\section{Introduction and Literature Review}

\subsection{Introduction}
The generation of terahertz radiation via the excitation of semiconductor surfaces with femtosecond laser pulses is a fundamental technique in modern photonics. This process relies on the acceleration of photo-generated carriers, where the resulting transient current acts as a Hertzian dipole source radiating electromagnetic waves into free space. While the physical mechanisms governing these currents—such as the Photo-Dember effect and surface depletion fields—are well-established in theoretical literature, the complex electrodynamics governing the subsequent propagation of the electric field often remain abstract and difficult to visualise.
\newline
The primary objective of this project is to bridge this gap by developing a dedicated simulation interface using MATLAB App Designer. Unlike standard numerical scripts, this graphical interface provides an interactive environment to model and visualise the generation of THz fields based on the method of Retarded Potentials. The specific motivation is to create a versatile tool that allows users to modify intrinsic semiconductor properties—such as carrier lifetime, diffusion coefficients, and absorption depth—to dynamically simulate the radiation characteristics of different materials. By transforming static theoretical equations into observable temporal waveforms and spatial vector maps, this interface serves as both a pedagogical tool for understanding surface emission physics and a predictive platform for optimising semiconductor materials in terahertz emitter design.

\subsection{Literature Review}
% TODO: Expand this into a proper literature review.
% - How is your project informed by existing knowledge?
% - How does existing knowledge relate to your project goals?
% - How will you use existing knowledge in your project?
% - How will you contribute to knowledge?
% Reference existing citations: \cite{Hebling:02}, \cite{SimulationTHz}, \cite{auston_picosecond_1984}

%% ============================================================
%% PROJECT GOALS
%% ============================================================
\section{Project Goals}
% TODO: Clearly state the goals/objectives of the project.
% - What are you trying to achieve?
% - What are the specific deliverables?

%% ============================================================
%% TECHNICAL BACKGROUND
%% ============================================================
\section{Technical Background}

\subsection{Photo-Dember Effect}
% TODO: Move the theoretical explanation of the Photo-Dember effect here.
% Present the background theory, mathematical framework, and physical principles.
The model is based on the diffusion-recombination theory: a highly non-equilibrium carrier distribution is generated in the semiconductor by ultra-short laser pulses and then evolves under the driving force of spatial gradients and mobility differences. Existing literature has pointed out that the time derivative of the instantaneous photogenerated current is the main source of terahertz radiation\cite{Hebling:02}. Therefore, the numerical framework of this project is based on this physical principle.

\subsection{Electric Field Theory}
The generation of terahertz radiation from ultrafast photocurrents is modelled using classical electrodynamics, where the source of the electromagnetic field is the time-varying current density, $\mathbf{J}(\mathbf{r}', t)$, distributed across the semiconductor surface. The radiated electric field $\mathbf{E}(\mathbf{r}, t)$ is derived using the method of retarded potentials and calculate use Jefimenko's equations. In the far-field approximation, where the observation distance $R$ is significantly larger than the source dimensions, the electrostatic and induction terms decay rapidly, leaving the radiative term as the dominant contribution. Consequently, the radiated electric field is directly proportional to the time derivative of the transverse current density\cite{SimulationTHz}.\newline
To account for the finite propagation speed of light $c$, the field observed at position $\mathbf{r}$ at time $t$ depends on the state of the current source at the retarded time $t_{ret} = t - R/c$. The polarization and directionality of the radiation are governed by the vector projection of the current acceleration onto the plane perpendicular to the line of sight $\hat{\mathbf{R}}$. This relationship is mathematically expressed as:
\begin{equation}
    \mathbf{E}(\mathbf{r}, t) \approx \frac{\mu_0}{4\pi} \iint_{S} \frac{\hat{\mathbf{R}} \times [\hat{\mathbf{R}} \times \dot{\mathbf{J}}(\mathbf{r}', t_{ret})]}{R} \, dS'
\end{equation}
where $\dot{\mathbf{J}}$ represents the partial time derivative of the current density. This formulation ensures that only the transverse component of the acceleration contributes to the propagating wave.

%% ============================================================
%% METHODS
%% ============================================================
\section{Methods}

\subsection{Photo-Dember Numerical Model}
First, a one-dimensional electron diffusion-composition model was constructed. By applying Gaussian excitation in both space and time, the expected carrier behavior was successfully simulated: rapid accumulation at the excitation center, symmetric diffusion over time, and exponential decay. The model output was consistent with the analytical expectations, verifying the stability and physical correctness of the numerical method.\newline
Subsequently, the model was extended to two dimensions to reflect the circular symmetry of the real light spot. The two-dimensional results showed that the carriers diffused radially symmetrically and an electron density peak slightly delayed relative to the excitation peak appeared, which was consistent with the ultrafast carrier dynamics reported in the literature. Further, by calculating the gradient of the electron density, the corresponding diffusion current was obtained. The result presented a typical transient current structure that radiated outward radially, which was in line with the characteristics of the Photo-Dember current.\newline
To enhance the physical accuracy of the model, additional holes were introduced and a double carrier coupling equation was constructed to reflect the mechanism of net current formation due to the difference in carrier mobilities between electrons and holes. The double carrier model shows that early transport is still dominated by electrons, while in the later stage, the spatial distribution of electrons and holes becomes more overlapping, resulting in a smoother overall carrier distribution, which is consistent with the existing experimental observations.

\subsection{Electric Field Computation}
The computational model discretises the semiconductor surface into a grid of elementary Hertzian dipoles. The calculation follows a three-step numerical algorithm. First, the time derivative of the current density, $\partial \mathbf{J}/\partial t$, is computed at each spatial grid point using a finite difference scheme. Second, since the calculated retarded time $t_{ret}$ does not generally align with discrete simulation steps, linear interpolation is employed to accurately estimate the current derivative at the exact retarded time. Finally, the principle of superposition is applied by summing the radiation contributions from all grid elements, weighted by the geometric decay factor $1/R$, to reconstruct the total temporal electric field waveform at the detector.

\subsection{MATLAB GUI Interface Design}
In this project, an integrated MATLAB App Designer GUI was designed and implemented to interface with the physics-based terahertz (THz) photocurrent simulation engine. The main goal is to provide a user-friendly, interactive platform that allows the user to configure material, excitation, and geometric parameters, execute the full 2D diffusion-recombination mode, and visualize the resulting carrier dynamics and THz waveforms.\newline
The overall GUI structure uses a modular, panel-based layout. Four functional tabs were implemented, including Mechanism Selector, Numeric Fields, Excitation, and $E(t)$ position. These tabs allow users to select the physical mechanism (built-in field, photo-Dember effect, or both), adjust semiconductor parameters (mobility, lifetime, absorption depth, temperature), configure optical excitation (pulse width, fluence, spot size, incidence angle, pulse shape), and specify the THz observation point in 3D space. A separate 3D axes window is implemented to visualize the rectangular semiconductor and the laser beam geometry, including features like incidence-angle-dependent beam rotation.\newline
Callback functions were developed that gather all GUI inputs, validate them, and package them into a structured parameter set for the simulation core. This includes converting fluence to peak intensity, mapping spot size to Gaussian waist, computing Einstein-relation diffusion coefficients, and applying incidence-angle correlations. The GUI supports additional physics toggles that enable the photo-Dember current and built-in drift current to be selectively activated.\newline
The GUI was integrated with the simulation backend, which returns the carrier density evolution, peak-density spatial distribution, and the THz electric field waveform computed via time-derivative current radiation. These results are visualized in three coordinated plots, including central carrier density vs. time, carrier distribution at the peak time, and $\lvert E \rvert (t)$ waveform at a specific chosen observation point.\newline
Consistent axis labels, physical units $(\mathrm{V/m},\ \mathrm{m^{-3}},\ \mu\mathrm{m})$, dynamic titles, and error-handling alerts were added to ensure robustness and scientific clarity.

%% ============================================================
%% RESULTS AND DISCUSSION
%% ============================================================
\section{Results and Discussion}

\subsection{Photo-Dember Model Results}
% TODO: Present and critically evaluate your Photo-Dember model results.
% Include figures, analysis, and comparison with literature.
Overall, the numerical platform has been established and verified, laying a solid foundation for subsequent calculations of terahertz radiation based on transient currents.
\begin{figure}[h]
	\begin{center}
	\includegraphics[width=14.5cm]{interim_report_latex_template (1)/figure/yuhao/image.png}
	\caption{
            Electron density heat map with vectors indicate current
	}
	\label{Electron density heat map with vectors indicate current}
	\end{center}
\end{figure}

\subsection{Electric Field Results}
The time-dependent behaviour of the radiated terahertz electric field was evaluated at a fixed far-field observation point located at $x = -20$ $\mu$m, $y = 30$ $\mu$m, $z = 10$ mm. Figure 2 presents the corresponding temporal traces of the $E_x$ and $E_y$ components. Both components display a distinctly bipolar pulse profile. This bipolar structure of the simulated field is in strong accordance with the theoretical model developed by Auston et al.\cite{auston_picosecond_1984}. In their pioneering study on picosecond photoconductors, Auston and co-workers showed that the far-field radiation produced by a transient current source is proportional to the first time derivative of the dipole moment, $\mathbf{E} \propto \dot{\mathbf{p}}(t)$. The results obtained here reproduce this anticipated behaviour: the leading positive part of the waveform is associated with the rapid rise of the photocurrent, whereas the following negative part is linked to the current decay. This confirms that the numerical implementation of the retarded potential accurately reflects the fundamental connection between carrier acceleration and electromagnetic radiation.
\begin{figure}[h]
	\begin{center}
	\includegraphics[width=9.5cm]{interim_report_latex_template (1)/figure/yaoqqi/field2d.png}
	\caption{
		Temporal waveforms of the radiated THz electric field components $E_x(t)$ (top) and $E_y(t)$ (bottom) calculated at an observation point of $(-20 \mu\text{m}, 30 \mu\text{m}, 10 \text{mm})$.
	}
	\label{field2d}
	\end{center}
\end{figure}

To characterise the polarisation of the emitted beam, the instantaneous transverse vector field distribution in the $xy$-plane was examined at a propagation distance of $z = 5$ mm and a time of $t = 0.2$ ps. As illustrated in Figure 3, the electric field vectors display clear radial symmetry, directed outward from the propagation axis.
\begin{figure}[h]
	\begin{center}
	\includegraphics[width=8.5cm]{interim_report_latex_template (1)/figure/yaoqqi/fielddistr.png}
	\caption{
		Instantaneous vector map of the THz electric field $(E_x, E_y)$ in the transverse plane at $z = 5$ mm. The field exhibits radial polarisation, consistent with the isotropic diffusion of photo-carriers from the excitation centre.
	}
	\label{fielddistr}
	\end{center}
\end{figure}

\subsection{GUI Interface Results}
% TODO: Present the GUI results with critical evaluation.
% How well does it work? Performance analysis?
\begin{figure}[h]
	\begin{center}
	\includegraphics[width=12.5cm]{interim_report_latex_template (1)/figure/shenglong/parameter.png}
	\caption{
		Parameters Selection Tabs
	}
	\label{Parameters Selection Tabs}
	\end{center}
\end{figure}

\begin{figure}[h]
	\begin{center}
	\includegraphics[width=12.5cm]{interim_report_latex_template (1)/figure/shenglong/3d.png}
	\caption{
		3D Semiconductor and Beam Visualization
	}
	\label{3D Semiconductor and Beam Visualization}
	\end{center}
\end{figure}

\begin{figure}[h]
	\begin{center}
	\includegraphics[width=12cm]{interim_report_latex_template (1)/figure/shenglong/central.png}
	\caption{
		Central Carrier Density Plot
	}
	\label{Central Carrier Density Plot}
	\end{center}
\end{figure}

\begin{figure}[b]
	\begin{center}
	\includegraphics[width=12cm]{interim_report_latex_template (1)/figure/shenglong/carrier.png}
	\caption{
		Carrier Distribution Plot at Peak Time
	}
	\label{Carrier Distribution Plot at Peak Time}
	\end{center}
\end{figure}
\begin{figure}[h]
	\begin{center}
	\includegraphics[width=12.5cm]{interim_report_latex_template (1)/figure/shenglong/Thz.png}
	\caption{
		THz Field at Specific Chosen Position and Incidence Angle
	}
	\label{THz Field at Specific Chosen Position and Incidence Angle}
	\end{center}
\end{figure}

\begin{figure}[h]
	\begin{center}
	\includegraphics[width=12.5cm]{interim_report_latex_template (1)/figure/shenglong/complete.png}
	\caption{
		Complete GUI Overview
	}
	\label{Complete GUI Overview}
	\end{center}
\end{figure}

%% ============================================================
%% CONCLUSIONS
%% ============================================================
\section{Conclusions}
% TODO: Write conclusions that summarise the main findings.
% - What is your main project result?
% - What did you achieve?
% - What lessons were learned?
% This replaces the old "Project management" and "Future work" sections.

%% ============================================================
%% TEAM CONTRIBUTIONS
%% ============================================================
\section{Team Contributions}
% TODO: List team member contributions and indicate who wrote which sections.
%
% Yaoqi Dong -- Photo-Dember effect modelling (1D, 2D, dual-carrier)
% Yuhao Chang -- Electric field computation (retarded potentials, 2D/3D field simulation)
% Shenglong Li -- MATLAB GUI design and integration

\newpage
\bibliography{example-references}
\bibliographystyle{IEEEtran} % IEEE style



\end{document}
